% Shift-uniform Combes--Thomas resolvents and covariance roots.
% Resolvent artifact: b542b3f2; positive-root artifact: bd11e55f.
\documentclass[11pt]{article}

\usepackage[margin=1.05in]{geometry}
\usepackage{amsmath,amssymb,amsthm,mathtools}
\usepackage{booktabs,longtable,array}
\usepackage{xcolor}
\usepackage[colorlinks=true,linkcolor=blue!55!black,
  urlcolor=blue!55!black,citecolor=blue!55!black]{hyperref}

\newtheorem{theorem}{Theorem}[section]
\newtheorem{lemma}[theorem]{Lemma}
\newtheorem{proposition}[theorem]{Proposition}
\newtheorem{corollary}[theorem]{Corollary}
\newtheorem{remark}[theorem]{Remark}
\theoremstyle{definition}
\newtheorem{definition}[theorem]{Definition}

\emergencystretch=3em
\tolerance=2000
\setlength{\parskip}{0.15em}
\raggedbottom

\newcommand{\R}{\mathbb{R}}
\newcommand{\N}{\mathbb{N}}
\newcommand{\SU}{\mathrm{SU}}
\newcommand{\Id}{\mathrm{Id}}
\newcommand{\dist}{\operatorname{dist}}
\newcommand{\op}{\mathrm{op}}
\newcommand{\inn}[2]{\langle #1,#2\rangle}
\newcommand{\norm}[1]{\lVert #1\rVert}
\newcommand{\lean}[1]{\texttt{#1}}
\newcommand{\formalized}{\textsc{Formalized}}
\newcommand{\proved}{\textsc{Proved}}
\newcommand{\openstatus}{\textsc{Open}}
\newcommand{\conditional}{\textsc{Conditional}}

\title{Shift-Uniform Combes--Thomas Resolvents\
on Physical Gauge Cochains:\
A Lean-Verified Route to Exponentially Localized Covariance Roots}

\author{Lluis Eriksson\thanks{This manuscript and part of the formal
development were produced with AI assistance under the repository's audit
protocol.  Claims labelled \formalized{} were checked by Lean~4 against a
pinned Mathlib and a targeted axiom oracle.  Paper-level functional-calculus
arguments are labelled \proved{} and are not represented as Lean theorems.}}
\date{August 1, 2026 -- research preprint, version 0.4}

\begin{document}
\maketitle

\begin{abstract}
Let $K$ be a coercive finite-range operator on the finite-dimensional real
Hilbert space of physical lattice gauge one-cochains.  A standard
Combes--Thomas argument controls the inverse kernel by conjugating $K$ with an
exponential weight.  Applying that argument naively to the shifted family
$K+s\Id$ appears to enlarge the entrywise bound by $s$, destroying the
uniformity needed to integrate resolvents.  We isolate and machine-check the
missing algebraic fact: scalar identity shifts cancel exactly from the tilt
defect,
\[
  (K+s\Id)_\theta-(K+s\Id)=K_\theta-K.
\]
Consequently, one fixed Combes--Thomas budget for $K$ controls every
$K+s\Id$, $s\ge0$.  If $K$ is coercive with constant $c>0$, has range $R$,
single-bond kernel bound $M$, and ball-cardinality bound $N_R$, then under
$M(e^{\theta R}-1)N_R\le c/2$, every right inverse $C_s$ of $K+s\Id$
satisfies
\[
  \norm{(C_s)_{q,p}}_{\op}
  \le \frac{2}{c+s}e^{-\theta\dist(q,p)}.
\]
The cancellation identity, shifted coercivity, uniform defect estimate, and
resolvent kernel theorem are formalized in Lean~4 on the repository's actual
physical gauge-cochain types.  At paper level, self-adjoint positivity and the
Stieltjes formula then give
\[
  \norm{(K^{-1/2})_{q,p}}_{\op}
  \le \frac{2}{\sqrt c}e^{-\theta\dist(q,p)}.
\]
The fractional-power decay itself is classical in matrix analysis; the
contribution claimed here is the shift-uniform, source-facing, machine-checked
route with explicit constants.  In version 0.4 the positive real continuous-
linear-map root is also constructed and verified spectrally in Lean, and a
new adapter derives every algebraic root-certificate field from covariance
positivity.  The remaining gap is stated without compression: the Stieltjes
identity and integrated single-bond root-kernel bound are not yet formalized,
so the localized physical root certificate remains conditional on that kernel
estimate; fixed-volume coercivity is still not volume-uniform.
\end{abstract}

\section{Claims and status discipline}\label{sec:claims}

This paper separates three levels that are easy to conflate.

\begin{center}
\begin{tabular}{p{0.20\textwidth}p{0.68\textwidth}}
\toprule
Status & Meaning in this manuscript \\
\midrule
\formalized & A named Lean theorem compiles on the pinned tree and the targeted
axiom oracle reports only \lean{propext}, \lean{Classical.choice}, and
\lean{Quot.sound}. \\
\proved & A complete finite-dimensional mathematical proof is given here, but
the proof is not yet a Lean declaration. \\
\openstatus & The statement is neither assumed silently nor claimed.  It is a
named frontier for subsequent formalization or analysis. \\
\bottomrule
\end{tabular}
\end{center}

The central resolvent result, Theorem~\ref{thm:shiftedCT}, is
\formalized.  The positive continuous-linear-map square root and its exact
square, positivity, and symmetry properties are now also \formalized.  The
Stieltjes passage to $K^{-1/2}$, Theorem~\ref{thm:root}, is \proved{} at paper
level.  Its identification with the formal spectral root, the integrated
single-bond kernel estimate, and unconditional discharge of the existing
\path{PhysicalLocalizedCovarianceRootCertificate} remain \openstatus.  No
continuum limit, mass gap, or volume-uniform positive coercivity constant is
claimed.

The mathematical phenomenon is not presented as new.  Exponential decay of
inverses of banded matrices goes back at least to Demko--Moss--Smith
\cite{DemkoMossSmith}; decay bounds for completely monotone functions,
including negative fractional powers, were developed directly by
Benzi--Simoncini \cite{BenziSimoncini}.  Our novelty claim is deliberately
narrower: a machine-checked shift-cancellation mechanism and a reusable
constant-explicit resolvent theorem on the physical gauge-cochain interface
used by the surrounding formal development.

\section{Physical cochains and localization predicates}\label{sec:setting}

Fix $d,N,N_c\in\N$ with $N>0$.  Let $\Lambda_N=(\mathbb Z/N\mathbb Z)^d$ and
let $\mathcal B_{d,N}$ be the finite set of positively oriented lattice bonds.
The internal coordinate space is the real coordinate model of
$\mathfrak{su}(N_c)$,
\[
  V_{N_c}\cong\R^{N_c^2-1}.
\]
The finite-dimensional real Hilbert space of physical one-cochains is
\[
  \mathcal H_{d,N,N_c}
  :=\ell^2(\mathcal B_{d,N};V_{N_c}).
\]
In Lean these are, respectively,
\lean{PhysicalBond d N}, \lean{SUNLieCoord Nc}, and
\lean{PhysicalGaugeOneCochain d N Nc}.

For $p\in\mathcal B_{d,N}$ and $v\in V_{N_c}$, write $\delta_pv$ for the
cochain supported at $p$ with value $v$.  This is
\lean{singlePhysicalBondCochain}.  Let
\[
  \dist:\mathcal B_{d,N}\times\mathcal B_{d,N}\longrightarrow\N
\]
be symmetric, satisfy the triangle inequality, and obey
$\dist(p,p)=0$.

Let $K:\mathcal H\to\mathcal H$ be a continuous linear map.  The formal
finite-range and kernel hypotheses have the following mathematical content:
\begin{align}
  \dist(q,p)>R &\Longrightarrow (K\delta_pv)(q)=0,
     \label{eq:range}\\
  \norm{(K\delta_pv)(q)}&\le M\norm v.
     \label{eq:kernel}
\end{align}
The cardinality parameter $N_R$ controls every metric ball,
\begin{equation}\label{eq:ball}
  \#\{y:\dist(x,y)\le R\}\le N_R.
\end{equation}
Finally, $K$ is coercive with constant $c$ if
\begin{equation}\label{eq:coercive}
  c\norm f^2\le\inn f{Kf}\qquad(f\in\mathcal H).
\end{equation}
The corresponding Lean predicate is \lean{IsCoerciveCLM K c}.

\section{Exponential conjugation}\label{sec:tilt}

Fix a root bond $r$ and $\theta\in\R$.  Define the diagonal tilt
\begin{equation}\label{eq:tilt}
  (T_{\theta,r}f)(q)=e^{\theta\dist(r,q)}f(q)
\end{equation}
and the conjugated operator
\begin{equation}\label{eq:conj}
  K_{\theta,r}=T_{\theta,r}KT_{-\theta,r}.
\end{equation}
Because the bond set is finite, these are bounded maps without a domain issue.
The formal development proves $T_{\theta,r}T_{-\theta,r}=\Id$ and, on a
single-bond probe,
\begin{equation}\label{eq:entry}
 (K_{\theta,r}\delta_pv)(q)
 =e^{\theta(\dist(r,q)-\dist(r,p))}(K\delta_pv)(q).
\end{equation}
On the range of $K$, the triangle inequality gives
$|\dist(r,q)-\dist(r,p)|\le R$.  The block Schur estimate in the formal
library then yields
\begin{equation}\label{eq:defect}
 \norm{K_{\theta,r}-K}
 \le \Delta_\theta,
 \qquad
 \Delta_\theta:=M(e^{\theta R}-1)N_R,
 \quad \theta\ge0.
\end{equation}
This is the existing Combes--Thomas engine.  The issue for a Stieltjes
construction is whether the same $\Delta_\theta$ controls all positive scalar
shifts.

\section{The scalar-shift cancellation}\label{sec:cancellation}

\begin{lemma}[identity is fixed; \formalized]\label{lem:id}
For every $r$, $\theta$, and physical gauge-cochain space,
\[
  T_{\theta,r}\Id T_{-\theta,r}=\Id.
\]
In Lean this is \lean{physicalTiltConjCLM\_id}.
\end{lemma}

\begin{proof}
It is the inverse relation $T_{\theta,r}T_{-\theta,r}=\Id$.  The Lean proof
unfolds the conjugation and rewrites with the already formalized composition
identity.
\end{proof}

\begin{lemma}[shift cancellation; \formalized]\label{lem:cancel}
For $s\in\R$,
\begin{equation}\label{eq:cancel}
  (K+s\Id)_{\theta,r}-(K+s\Id)=K_{\theta,r}-K.
\end{equation}
In Lean this is
\lean{physicalTiltConjCLM\_add\_smul\_id\_sub}.
\end{lemma}

\begin{proof}
Conjugation is linear in the operator.  Lemma~\ref{lem:id} gives
$(K+s\Id)_{\theta,r}=K_{\theta,r}+s\Id$; subtracting $K+s\Id$ proves
\eqref{eq:cancel}.  The Lean proof separately establishes additivity and
scalar compatibility of conjugation and closes the operator identity by
abelian-group normalization.
\end{proof}

The importance of \eqref{eq:cancel} is not its algebraic difficulty but its
location in the estimate.  A naive application of the unshifted theorem to
$K+s\Id$ asks for a kernel bound on the whole shifted operator and therefore
introduces a spurious term proportional to $s$.  Equation~\eqref{eq:cancel}
shows that the tilt only sees the off-diagonal commutator; the scalar diagonal
shift is invisible.

\begin{corollary}[uniform defect; \formalized]\label{cor:defect}
Under \eqref{eq:range}--\eqref{eq:ball}, for every $s\in\R$,
\[
 \norm{(K+s\Id)_{\theta,r}-(K+s\Id)}\le\Delta_\theta.
\]
This is \lean{norm\_physicalTiltConj\_add\_smul\_id\_sub\_le}.
\end{corollary}

\section{Shift-uniform coercivity and resolvents}\label{sec:resolvent}

The scalar shift also improves coercivity exactly.

\begin{lemma}[shifted coercivity; \formalized]\label{lem:coershift}
If \eqref{eq:coercive} holds, then
\[
  (c+s)\norm f^2\le\inn f{(K+s\Id)f}
\]
for every real $s$.  In Lean this is
\lean{isCoerciveCLM\_add\_smul\_id}.
\end{lemma}

Combining Lemma~\ref{lem:coershift} with
Corollary~\ref{cor:defect} gives the shift-uniform CT2 endpoint
\begin{equation}\label{eq:tiltcoer}
 \inn f{(K+s\Id)_{\theta,r}f}
 \ge(c+s-\Delta_\theta)\norm f^2.
\end{equation}
This is
\lean{isCoerciveCLM\_physicalTiltConj\_add\_smul\_id}.
If $s\ge0$ and $\Delta_\theta\le c/2$, then
\begin{equation}\label{eq:half}
  c+s-\Delta_\theta\ge\frac{c+s}{2}.
\end{equation}
The corresponding formal theorem is
\lean{isCoerciveCLM\_physicalTiltConj\_add\_smul\_id\_half}.

\begin{theorem}[shift-uniform Combes--Thomas resolvent; \formalized]
\label{thm:shiftedCT}
Assume $c>0$, $M\ge0$, $\theta>0$, equations
\eqref{eq:range}--\eqref{eq:coercive}, and
\begin{equation}\label{eq:budget}
  M(e^{\theta R}-1)N_R\le\frac c2.
\end{equation}
For $s\ge0$, let $C_s$ be any right inverse of $K+s\Id$:
\[
  (K+s\Id)C_s=\Id.
\]
Then, for all bonds $p,q$ and $v\in V_{N_c}$,
\begin{equation}\label{eq:CTs}
  \norm{(C_s\delta_pv)(q)}
  \le \frac{2}{c+s}
       e^{-\theta\dist(q,p)}\norm v.
\end{equation}
This is
\lean{physicalCovariance\_exponentialKernelBound\_of\_coercive\_add\_smul\_id}.
\end{theorem}

\begin{proof}[Proof architecture]
Root the tilt at the source bond $p$ and set
$y=T_{\theta,p}C_s\delta_pv$.  Since $\dist(p,p)=0$, the source probe is fixed
by its own tilt.  The right-inverse relation and exact untilting give
\[
  (K+s\Id)_{\theta,p}y=\delta_pv.
\]
Coercivity at $(c+s)/2$ and Cauchy--Schwarz imply
\[
  \norm y\le\frac{2}{c+s}\norm v.
\]
At the target bond,
\[
  (C_s\delta_pv)(q)=e^{-\theta\dist(p,q)}y(q).
\]
The coordinate projection has norm at most one, and symmetry of the distance
gives \eqref{eq:CTs}.  Every equality and inequality in this proof architecture
is represented in the cited Lean theorem.
\end{proof}

\begin{remark}[why the denominator is integrable]
The key improvement over a shift-blind estimate is not merely that the decay
rate remains positive.  The amplitude retains the spectral shift as
$(c+s)^{-1}$.  Substituting $s=t^2$ produces an integrable majorant on
$[0,\infty)$, which is exactly what the inverse-square-root Stieltjes formula
requires.
\end{remark}

\section{The covariance-root corollary}\label{sec:root}

The next theorem is ordinary finite-dimensional spectral calculus.  It is
included with a complete proof to identify the exact remaining Lean task.

\begin{theorem}[localized inverse square root; \proved]\label{thm:root}
In addition to the hypotheses of Theorem~\ref{thm:shiftedCT}, assume that $K$
is self-adjoint and positive with
$\sigma(K)\subset[c,\infty)$, $c>0$.  Then the unique positive inverse square
root $K^{-1/2}$ satisfies
\begin{equation}\label{eq:rootbound}
  \norm{(K^{-1/2}\delta_pv)(q)}
  \le\frac{2}{\sqrt c}
      e^{-\theta\dist(q,p)}\norm v.
\end{equation}
Moreover,
\[
  \norm{K^{-1/2}}\le c^{-1/2},\qquad
  K^{-1/2}K^{-1/2}=K^{-1}.
\]
\end{theorem}

\begin{proof}
The spectral theorem gives the norm-convergent Stieltjes identity
\begin{equation}\label{eq:stieltjes}
  K^{-1/2}=\frac{2}{\pi}\int_0^\infty(K+t^2\Id)^{-1}\,dt.
\end{equation}
Indeed, after diagonalizing $K$, the operator identity reduces to the scalar
identity, valid for $\lambda>0$,
\begin{equation}\label{eq:scalar}
 \frac{2}{\pi}\int_0^\infty\frac{dt}{\lambda+t^2}
 =\frac{2}{\pi}\frac{\pi}{2\sqrt\lambda}
 =\lambda^{-1/2}.
\end{equation}
There is no interchange difficulty in finite dimension.  Apply
Theorem~\ref{thm:shiftedCT} with $s=t^2$ and integrate the pointwise norm
majorant:
\begin{align*}
 \norm{(K^{-1/2}\delta_pv)(q)}
 &\le \frac{2}{\pi}\int_0^\infty
       \frac{2}{c+t^2}e^{-\theta\dist(q,p)}\norm v\,dt\\
 &=\frac{2}{\sqrt c}e^{-\theta\dist(q,p)}\norm v.
\end{align*}
The spectral bounds and the square identity follow by applying
$\lambda\mapsto\lambda^{-1/2}$ on $\sigma(K)$.
\end{proof}

\begin{remark}[constant accounting]
The spectral norm of $K^{-1/2}$ is at most $c^{-1/2}$, whereas the kernel
amplitude obtained from the half-budget CT argument is $2c^{-1/2}$.  The
factor two is the price of reserving half the coercivity uniformly for every
tilt.  We make no optimality claim for this constant or for the admissible
rate $\theta$.
\end{remark}

\section{Connection to the formal covariance-root frontier}\label{sec:frontier}

The repository already contains the proposition-valued structure
\texttt{PhysicalLocalizedCovariance}\allowbreak\texttt{RootCertificate}.
Given a precision $K$, a
covariance $C$, and a proposed root $B$, it requires:
\begin{enumerate}
\item a localized covariance certificate for $(K,C)$;
\item $B\circ B=C$;
\item a norm bound on $B$;
\item self-adjointness of $B$ in inner-product form;
\item positivity of $B$;
\item a single-bond kernel bound on $B$.
\end{enumerate}
Theorem~\ref{thm:root} supplies the mathematical content of items 2--6 when
$C=K^{-1}$ and $B=K^{-1/2}$, while Theorem~\ref{thm:shiftedCT} supplies the
uniform resolvent estimate needed for item 6.  The new module
\path{FiniteDimensionalRealPositiveSqrt.lean} now supplies a Lean term $B$ for
every positive finite-dimensional real continuous endomorphism and proves
$B\circ B=C$, positivity, and symmetry.  The physical adapter derives items
2--5 (using $\norm B$ itself for the non-quantitative norm bound); item 6 is
still an explicit analytic input.

The local-SPD route in
\lean{PhysicalGaugeLocalSPDPrecisionRoot.lean} is complementary.  It develops
Neumann and inverse-square-root coefficient majorants under a normalized
small-perturbation hypothesis
$K=a(\Id-L)$ with $\norm L<1$, but deliberately leaves the continuous-linear
root and its kernel bound as source fields.  The present route replaces that
smallness hypothesis by coercivity plus finite range.  Version 0.4 closes the
positive-root object for every symmetric coercive precision; neither route yet
proves the physical single-bond kernel bound for that canonical root.

The exact next formal chain is therefore
\begin{equation}\label{eq:nextchain}
\boxed{
\begin{gathered}
\text{self-adjoint coercive }K
\longrightarrow \text{formal positive spectral root }B,\quad B^2=K^{-1}
\\
\longrightarrow \text{formal identification with the Stieltjes integral}
\longrightarrow \text{integrated single-bond kernel bound}
\\
\longrightarrow
\lean{PhysicalLocalizedCovarianceRootCertificate}.
\end{gathered}}
\end{equation}
The first arrow is now machine-checked by transporting a real positive-
semidefinite matrix root through an orthonormal basis.  The remaining
engineering question is to identify this spectral root with the
norm-convergent Stieltjes integral while preserving the single-bond kernel API.

\paragraph{Post-review implementation audit.}
The available library route was not a one-line application of an existing
operator square-root theorem.  The implementation chooses an orthonormal
coordinate model, applies the real matrix spectral theorem, takes the
nonnegative square root of every eigenvalue, and transports the result back to
a continuous linear map.  Lean verifies positive semidefiniteness and the
exact square before transport, then verifies the continuous-linear-map
properties after transport.  This closes the first two obligations identified
in version 0.3.1.  The two remaining obligations are the norm-convergent
Stieltjes identification and the integrated single-bond kernel estimate.
Supplying the last predicate as an assumption still does not prove decay; the
new certificate adapter uses it only to isolate the exact remaining frontier.

\section{Prior art and exact delta}\label{sec:prior}

Combes and Thomas introduced exponential conjugation in spectral analysis
\cite{CombesThomas}.  Demko, Moss, and Smith proved exponential decay of
inverse band matrices with spectral-rate information
\cite{DemkoMossSmith}.  Benzi and Simoncini treated completely monotone
functions of Hermitian banded and sparse matrices, explicitly including
negative fractional powers \cite{BenziSimoncini}.  Thus neither
Theorem~\ref{thm:root} nor the use of a resolvent integral is claimed as a new
analytic phenomenon.

\begin{center}
\small
\begin{tabular}{p{0.19\textwidth}p{0.23\textwidth}p{0.46\textwidth}}
\toprule
Antecedent & Overlap & Exact difference claimed here \\
\midrule
Combes--Thomas \cite{CombesThomas} & Exponential weighting and resolvent
localization & Finite physical gauge-cochain API, explicit Schur budget, exact
scalar-shift cancellation, and machine-checked source-rooted extraction. \\
Demko--Moss--Smith \cite{DemkoMossSmith} & Off-diagonal decay of inverses &
Metric finite-range block operators rather than scalar band matrices; the
main delta is formal verification, not a sharper rate. \\
Benzi--Simoncini \cite{BenziSimoncini} & Decay of negative fractional powers
via matrix-function representations & No claim of analytic priority.  We
provide a Lean-verified resolvent family with the precise $(c+s)^{-1}$
amplitude consumed by the physical covariance-root interface. \\
Repository fixed-volume CT chain \cite{ErikssonPoincare} & Formal inverse
decay with amplitude $2/c$ & The old theorem did not preserve a scalar shift;
the new theorem is uniform for all $s\ge0$ and enables the Stieltjes integral. \\
Repository local-SPD root frontier & Inverse-square-root coefficient and kernel
majorants & Present hypotheses are coercivity plus finite range, not
$\norm L<1$; the canonical real root is now formal, while its localized kernel
bound remains open. \\
\bottomrule
\end{tabular}
\end{center}

The literature search for this manuscript was performed on August 1, 2026,
using combinations of ``Combes--Thomas'', ``inverse square root'',
``fractional powers'', ``sparse matrix functions'', ``Stieltjes'', and
``Lean''.  We found no machine-checked Combes--Thomas inverse-square-root
theorem, but this negative search result is not an exhaustiveness claim.  The
paper's defensible dominance axis is reproducibility and integration into a
concrete formal physics interface, not classical analytic novelty.

\section{Adversarial audit and limitations}\label{sec:audit}

\paragraph{Attack 1: the theorem is already known under matrix-function
language.}  Correct at the analytic level.  This is why the title says
``Lean-verified route'' and not ``new localization theorem'', and why
Section~\ref{sec:prior} treats negative fractional powers as direct prior art.

\paragraph{Attack 2: the result is only finite volume.}  Correct.  The formal
physical cochain space is finite-dimensional.  A separate machine-checked
Poincar\'e-wall result shows that the current unscaled block-Poincar\'e route
does not provide a positive coercivity constant uniform in volume for the
relevant dimensions \cite{ErikssonPoincare}.  The present theorem preserves
whatever $c$ is supplied; it does not manufacture a uniform $c$.

\paragraph{Attack 3: coercivity does not imply a positive square root for a
non-self-adjoint operator.}  Correct.  The formal resolvent theorem does not
need self-adjointness, but Theorem~\ref{thm:root} adds self-adjoint positivity
explicitly.  Omitting it would be a false bridge.

\paragraph{Attack 4: a right inverse family is not a measurable or continuous
resolvent map.}  In finite dimension the canonical inverse
$(K+t^2\Id)^{-1}$ is continuous in $t$ for a positive self-adjoint $K$.  The
Lean theorem is intentionally pointwise in an arbitrary right inverse; the
continuity and its integral identification with the now-formal spectral root
belong to the open Stieltjes layer.

\paragraph{Attack 5: the diagonal is not genuinely off-diagonal data.}  The
formal global kernel bound on the base operator still includes diagonal
entries.  The improvement is that the \emph{shift} is eliminated exactly from
the tilt defect.  We do not claim a new off-diagonal norm predicate in this
version.

\paragraph{Attack 6: the constants may be poor.}  Correct.  The factor two and
the Schur cardinality $N_R$ are explicit but not optimized.  Polynomial or
best-approximation methods may give sharper rates for banded Hermitian
matrices.  The purpose here is a robust formal interface.

\paragraph{Kill criteria.}  The programme should be reformulated if any of the
following occurs: (i) the exact shift identity fails under the actual root
representation; (ii) the functional-calculus bridge cannot preserve the
single-bond kernel API without assuming the desired kernel bound; (iii) the
root certificate requires a volume-uniform $c$ from the already refuted block
normalization; or (iv) a prior formalization is found with the same physical
types, hypotheses, constants, and endpoint.  None is currently observed, but
(ii) is the real technical bottleneck.

\section{Lean theorem map and reproducibility}\label{sec:repro}

\subsection{Theorem-to-artifact map}

The resolvent source modules are
\path{YangMills/RG/PhysicalCoerciveCombesThomas.lean} (CT1/CT2) and
\path{YangMills/RG/PhysicalCoerciveCombesThomasInverse.lean} (CT3/CT4).
The root layer is implemented in
\path{FiniteDimensionalRealPositiveSqrt.lean},
\path{CoerciveCovariancePositiveSqrt.lean}, and
\path{PhysicalGaugeCovariancePositiveRoot.lean}.

\footnotesize
\begin{longtable}{@{}p{0.41\textwidth}p{0.32\textwidth}p{0.13\textwidth}@{}}
\toprule
Lean declaration & File & Status \\
\midrule
\endhead
\path{physicalTiltConjCLM_id} &
CT1/CT2 module & \textsc{exact} \\
\path{physicalTiltConjCLM_add_smul_id_sub} & CT1/CT2 module & \textsc{exact} \\
\path{isCoerciveCLM_add_smul_id} & CT1/CT2 module & \textsc{exact} \\
\path{norm_physicalTiltConj_add_smul_id_sub_le} & CT1/CT2 module & \textsc{exact} \\
\path{isCoerciveCLM_physicalTiltConj_add_smul_id} & CT1/CT2 module & \textsc{exact} \\
\path{isCoerciveCLM_physicalTiltConj_add_smul_id_half} &
CT3/CT4 module & \textsc{exact} \\
\path{physicalCovariance_exponentialKernelBound_of_coercive_add_smul_id}
& same & \textsc{exact} \\
\path{finiteDimensionalRealPositiveSqrt_}\newline
\path{comp_self} &
\path{FiniteDimensionalRealPositive}\newline\path{Sqrt.lean} & \textsc{exact} \\
\path{finiteDimensionalRealPositiveSqrt_}\newline
\path{isPositive} & same & \textsc{exact} \\
\path{covarianceSqrtOfIsCoerciveCLM_comp_self} &
\path{CoerciveCovariancePositive}\newline\path{Sqrt.lean} & \textsc{exact} \\
\path{physicalLocalizedCovariance}\newline
\path{RootCertificate_of_positive_covariance} &
\path{PhysicalGaugeCovariance}\newline\path{PositiveRoot.lean} & \textsc{conditional} \\
Stieltjes identity and Theorem~\ref{thm:root} & this manuscript & paper \\
Integrated root-kernel bound & not yet implemented & open \\
\bottomrule
\end{longtable}
\normalsize

\subsection{Snapshot and commands}

Repository:
\url{https://github.com/lluiseriksson/THE-ERIKSSON-PROGRAMME}.
The shift-uniform resolvent artifact is the dedicated Git commit on branch
\lean{codex/stieltjes-root}
\begin{center}
\small\lean{b542b3f2f066b60a914977f2a50160e403a53367}.
\end{center}
Its parent is \lean{81e9cf064d6771f986978ccffb96018c78b5e459}.  The commit contains exactly
the two theorem modules, the oracle registrations, and the verification-ledger
entry.  Thus the formal source cited here is an immutable Git object rather
than an uncommitted working-tree patch.  Remote publication and a hosted CI run
of that commit remain release actions and are not claimed by this manuscript.
The positive-root extension is the later immutable commit
\begin{center}
\small\lean{bd11e55f9cfe572d5ec4200f49c45339df4e3c03}.
\end{center}
It adds the three root
modules, their aggregate imports, oracle registrations, and verification-ledger
Addendum 571.  The manuscript edits follow that commit, so the root source hash
is not circular.
The version-0.2 manuscript artifact is commit
\lean{25f532eaeecbfb66ee146e264953cde0e2e0ea29}.  Version 0.3 is distributed
under the annotated local release tag \lean{shift-uniform-ct-v0.3} together
with a complete Git source archive, rather than as a PDF alone.  The archive
contains this manuscript, the pinned \path{lake-manifest.json}, both Lean
modules, \path{oracle_check.lean}, the focused
\path{ShiftUniformOracle.lean}, and a verification transcript.  A reader can
therefore materialize and check the cited source without relying on an
uncommitted worktree.

For cryptographic continuity, the materialized version-0.3 source package is
the tree at commit
\lean{cba2b5d705bfadc377a9c04293a57a6bfcb82c19}.  Its exact file name is
\path{shift-uniform-combes-thomas-v0.3-source.zip}; it contains 7,350 archive
entries and has SHA-256
\begin{center}
\small\lean{1B5EAB37038E071F1E3A1A1BA263D93048828EA079950273BF915B4605897823}.
\end{center}
Version 0.3.1 was the documentation-only revision that printed those package
identifiers in the PDF.  Version 0.4 is a scientific revision: it cites the
separate positive-root commit above and updates the formal/open boundary.  It
does not alter the immutable version-0.3 source archive or its hash.

Toolchain: Lean \path{leanprover/lean4:v4.29.0-rc6}; Mathlib revision
\texttt{0764272048}\ldots{} (the full revision is recorded in
\path{lake-manifest.json}).  The focused checks are:
\begin{verbatim}
lake env lean YangMills/RG/PhysicalCoerciveCombesThomas.lean
lake env lean YangMills/RG/PhysicalCoerciveCombesThomasInverse.lean
lake build +YangMills.RG.PhysicalCoerciveCombesThomas:olean
lake build +YangMills.RG.PhysicalCoerciveCombesThomasInverse:olean
lake build YangMills.RG.FiniteDimensionalRealPositiveSqrt
lake build YangMills.RG.CoerciveCovariancePositiveSqrt
lake build YangMills.RG.PhysicalGaugeCovariancePositiveRoot
\end{verbatim}
The four resolvent checks and all three root-module builds completed
successfully on August 1, 2026.  Two additional attempts to rebuild the full
\path{YangMillsCore} aggregator were terminated by the command wrapper after
10 and 20 minutes, respectively, without a Lean error; no successful fresh
full-core rebuild is claimed.  A targeted oracle checked the seven resolvent
declarations plus five root declarations; each printed exactly
\[
 [\lean{propext},\ \lean{Classical.choice},\ \lean{Quot.sound}].
\]
The declarations are also registered in \path{oracle_check.lean}.  The
focused audit is independently runnable as
\begin{verbatim}
lake env lean papers/shift-uniform-combes-thomas/ShiftUniformOracle.lean
\end{verbatim}
The development contains no \lean{sorry} in the two resolvent modules or the
three new root modules.

\paragraph{External author-corpus provenance note.}
As of August 1, 2026, thirteen corrective or supersession forms concerning
other manuscripts had been sent and remained \textbf{ENVIADO/PENDIENTE}; none
should be resent or described as published while the public records lag.
The scope split between 2607.0035 and 2607.0039 remains controlling, and
2607.0089 remains \textbf{REVIEW-PENDING}.  The complete record-level matrix,
promotion rule, and checks are moved to
\path{papers/shift-uniform-combes-thomas/PUBLICATION-PROVENANCE.md}.  They are
administrative provenance, not hypotheses, evidence, or citations for any
mathematical claim in this paper.

\section{Conclusion}\label{sec:conclusion}

The useful observation is exact and small:
scalar identity shifts disappear from exponential-conjugation defects.  Its
consequence is structurally larger: a single fixed Combes--Thomas budget
controls an entire positive resolvent family with the integrable amplitude
$(c+s)^{-1}$.  That family is now a machine-checked theorem on the physical
gauge-cochain types that the covariance-root interface actually consumes.

The paper-level Stieltjes integral then produces exponential localization of
$K^{-1/2}$ with amplitude $2/\sqrt c$.  This is not advertised as a new result
in matrix analysis.  Version 0.4 does close one previously open bridge: Lean
now constructs the positive real continuous-linear root and derives its exact
square, positivity, symmetry, and the corresponding algebraic certificate
fields.  What remains is narrower and genuinely analytic: identify that root
with the Stieltjes integral and transfer the integrated resolvent estimate to
the single-bond root-kernel predicate.  Until that theorem is formalized, the
localized physical certificate remains conditional; the fixed-volume
Poincar\'e wall is unchanged.

\begin{thebibliography}{99}

\bibitem{CombesThomas}
J.-M.~Combes and L.~Thomas,
\emph{Asymptotic behaviour of eigenfunctions for multiparticle
Schr\"odinger operators},
Communications in Mathematical Physics \textbf{34} (1973), 251--270.
\href{https://doi.org/10.1007/BF01646473}{doi:10.1007/BF01646473}.

\bibitem{DemkoMossSmith}
S.~Demko, W.~F.~Moss, and P.~W.~Smith,
\emph{Decay rates for inverses of band matrices},
Mathematics of Computation \textbf{43} (1984), 491--499.
\href{https://doi.org/10.1090/S0025-5718-1984-0758197-9}
{doi:10.1090/S0025-5718-1984-0758197-9}.

\bibitem{BenziSimoncini}
M.~Benzi and V.~Simoncini,
\emph{Decay bounds for functions of Hermitian matrices with banded or
Kronecker structure},
SIAM Journal on Matrix Analysis and Applications \textbf{36} (2015),
1263--1282.
\href{https://doi.org/10.1137/151006159}{doi:10.1137/151006159}.

\bibitem{BalabanI}
T.~Ba{\l}aban,
\emph{Propagators and renormalization transformations for lattice gauge
theories. I},
Communications in Mathematical Physics \textbf{95} (1984), 17--40.

\bibitem{BalabanII}
T.~Ba{\l}aban,
\emph{Propagators and renormalization transformations for lattice gauge
theories. II},
Communications in Mathematical Physics \textbf{96} (1984), 223--250.

\bibitem{Lean4}
L.~de~Moura and S.~Ullrich,
\emph{The Lean 4 theorem prover and programming language},
CADE-28, Lecture Notes in Computer Science 12699 (2021), 625--635.

\bibitem{mathlib}
The mathlib community,
\emph{The Lean mathematical library},
CPP 2020, 367--381.

\bibitem{ErikssonPoincare}
L.~Eriksson,
\emph{The volume-uniform Poincar\'e walls: machine-checked obstructions for
flat and fluctuation-sector block-Poincar\'e routes to Combes--Thomas
coercivity in lattice Yang--Mills},
repository preprint (2026),
\path{papers/poincare-wall/poincare_wall.tex}.

\end{thebibliography}

\end{document}
